\documentclass[11pt]{article}
\usepackage[bottom]{footmisc} 
\renewcommand{\thempfootnote}{\arabic{mpfootnote}}

\usepackage[hidelinks]{hyperref}

\usepackage[utf8]{inputenc}

\usepackage[T1]{fontenc}
\usepackage[ttdefault=true]{AnonymousPro}

\usepackage{multicol}
\setlength{\columnsep}{1cm}

\usepackage{abc}
\usepackage{graphicx}
\usepackage{musicography}

\usepackage[margin=0.5in]{geometry} 
\geometry{letterpaper}


\usepackage{fancyhdr} 
\pagestyle{fancy} % options: empty , plain , fancy
\renewcommand{\headrulewidth}{0pt} % customise the layout...
\lhead{}\chead{}\rhead{}
\lfoot{}\cfoot{\thepage}\rfoot{}


\begin{document}

\thispagestyle{empty}

\begin{abc}[name=durangs,options={-O= +c -F single.fmt}]
X:1
T:Durang's Hornpipe
T:Wobble Gears
C:Howard Marshall's version
C:transcribed by Kitty Steetle
M:C|
K:D
FE |: "D"DDFA DDFA | "G"BcdA BAFE | "D"DEFD EDFD | "A"E2 CC A,CCE | 
"D"DDFA DDFA | "G"BcdA BA^GB | "A"ABce fecA |1 "D"dBAF D2FE :|2 "D"dBAF D4 |:
"D"fgae fedA | "G"BcdA BAFD | "D"G4 F4 | "A"E4 A4 | 
"D"fgae fedA | "G"BcdA BA^GB | "A"ABce fecA | "D"dBAF D4 :|
\end{abc}

\vskip 0.5in
\noindent\makebox[\textwidth][c]{%
  \begin{minipage}{7.0in}
    \setlength{\parindent}{1.5em}
    
The tune, Durang's Hornpipe, is believed to be composed in 1785 by Mr
Hoffmaster and his wife, both little people, and named for their
fiddle student and actor/dancer, John Durang (Lancaster, PA
1768 -- Philadelphia, PA 1821). John Durang then went on to be known as
`the first American dancer' and popularized the `nautical-style'
hornpipe dance, The Sailor's
Hornpipe.\footnote{\url{https://tunearch.org/wiki/Annotation:Durang\%27s\_Hornpipe\_(1)}
  This sailor's hornpipe dance is for a single dancer or group of
  single dancers as opposed to the hornpipe set dance for multiple
  couples.}

By the early 1800s the tune is found in American music commonplace
books\footnote{Commonplace books are an `individual's compilation of
  memorabilia (proverbs, poems, recipes, or on occasion pieces of
  music) recorded by hand in blank books'. (Goertzen, \textbf{American
    Antebellum Fiddling}, p 1)} including Thomas W Shannon, thought to
be a teacher of Arthur McArthur Sr (Scotland, 1790 -- ME,
1874).\footnote{Goertzen, \textit{op. cit.}, p 4.} It also appears in
commonplace books of Ebenezer Bevens (Middletown, CT, 1825), Gurden
Trumbull (Stonington, CT, 1801), Josiah Adams (Framingham, MA, 1808),
and Seth Johnson (Woburn, MA, 1807).\footnote{tunearch,
  \textit{loc. cit.}}  Durang's Hornpipe is a popular standard tune in
Old Time Jams and Irish Sessions throughout North America.

---Kitty Steetle

\textit{Josh adds: there are many versions of this tune. Since this
  isn't one I play a lot, the Tune of the Month recording is most
  likely a blend of the last few versions I'd listened to.}

  \end{minipage}
  \vspace*{\fill}
}

\end{document}
